In modern society, we are constantly surrounded by environments designed to capture our attention through potent sensory enticements—such as bright lights, vivid colors, and distinctive sounds. From smartphone notifications to dynamic digital advertisements, these attention-grabbing stimuli have become ubiquitous in our daily lives. Perhaps nowhere is this phenomenon more deliberately engineered than in gambling environments, particularly in electronic gambling machines (EGMs). When a player wins on a digital slot machine, they are greeted not merely with monetary reward, but with an orchestrated sensory celebration—flashing lights, triumphant sounds, and animated displays—creating a multisensory experience far more salient than the simple fact of winning money. This deliberate design strategy is not merely aesthetic; research indicates that these audiovisual features may influence gambling persistence—especially for higher‑risk gamblers—with players demonstrating longer play sessions on machines that provide enhanced sensory feedback.

This strategic pairing of rewards with intensified sensory feedback raises a fundamental question about human cognition: How do salient reward-paired cues (SRC) influence the way we learn from outcomes and make subsequent decisions? The answer to this question has implications extending far beyond gambling contexts—it touches upon basic mechanisms of learning, decision-making, and ultimately behaviors that can range from adaptive to potentially harmful. From a clinical perspective, understanding these mechanisms may provide insights into vulnerability factors for conditions like gambling disorder, which affects approximately 1.2\% of the population and is associated with significant psychological, social, and financial harm \citep{who2024}. Therefore, this thesis investigates how such SRCs, inspired by EGM design, impact learning dynamics and choice behavior in a controlled setting.

\subsection{Reinforcement Learning as a Framework for Understanding Decision-Making}

Reinforcement learning (RL) provides a powerful and computationally precise framework for understanding how humans and other animals learn from the consequences of their actions. At its core, RL posits that agents adjust their expectations and behaviours based on prediction errors—the discrepancy between expected and received outcomes \citep{sutton2020}. When we experience an unexpected reward, a positive prediction error signal drives us to increase the value we assign to the action that led to that reward, making us more likely to repeat it in the future. Conversely, unexpected negative outcomes lead us to devalue certain choices and avoid them subsequently.

A critical parameter in reinforcement learning models is the learning rate (alpha, $\alpha$), which determines how quickly value estimates are updated in response to new information. Formally, the value update can be expressed as: 
\begin{equation*}
V_{\text{new}} = V_{\text{old}} + \alpha \times (R - V_{\text{old}})
\end{equation*}
where $V$ represents the value estimate before ($V_{\text{old}}$) and after ($V_{\text{new}}$) receiving the outcome, $R$ is the outcome received, $\alpha$ represents the learning rate, typically constrained between 0 and 1 \citep{sutton2020}. The term $(R-V_{\text{old}})$ quantifies the prediction error, representing the difference between the actual reward obtained and the value that was expected. Higher learning rates result in more rapid updating based on recent outcomes, while lower rates produce more stable value estimates that change gradually over time. Importantly, research suggests that learning rates are not fixed but can be dynamically modulated by contextual factors, emotional states, and attentional processes \citep{lepelley2016, pike2022, wu2017}. This adaptive flexibility in learning rates may be advantageous in volatile environments but could potentially be exploited by artificial environments designed to maximize engagement.

Beyond these formal value-updating mechanisms, decision-making can sometimes be guided by simpler heuristics. One such fundamental strategy is ``win-stay, lose-shift'' (WSLS). This heuristic dictates a straightforward policy: if a chosen action leads to a positive outcome (a win), repeat that action on the next opportunity; if it leads to a negative outcome (a loss), switch to an alternative action \citep{robbins1952, nowak1993}. While computationally less demanding than maintaining and updating precise value estimates, WSLS can be surprisingly effective, particularly in environments where outcomes are somewhat stable. In the context of the current study, examining the win-stay component of this strategy provides a direct behavioral index of how recent rewards influence subsequent choices. It is therefore pertinent to investigate whether factors like feedback salience modulate not only the computationally derived learning rate ($\alpha$) but also this more basic tendency to repeat rewarded actions.

Traditional RL models often focus on the valence (positive/negative) and magnitude of outcomes, potentially overlooking how the sensory presentation of these outcomes—their salience—might directly influence the learning process. Salience refers to the degree to which a stimulus stands out and captures attention, which can be due to its sensory intensity (e.g., brightness, loudness), novelty, or motivational significance \citep{berridge1998, itti2001}. Empirical studies have demonstrated that manipulating the visual salience of stimuli can significantly bias attention allocation and subsequent decision-making \citep{towal2013}, and that salient perceptual features can override value-based decision processes under certain conditions \citep{tsetsos2012}.

\subsection{Electronic Gambling Machines and Psychological Salience}

The gambling industry appears to intuitively understand the power of salient feedback, as evidenced by the sophisticated sensory design of modern EGMs. These machines create an immersive environment where wins—and notably, ``losses disguised as wins'' (LDWs)—are celebrated with disproportionate audiovisual feedback. LDWs occur when outcomes that are net losses (smaller than the initial wager) are nevertheless accompanied by celebratory sensory stimuli typically associated with wins, confusing players about their actual financial outcomes and contributing to an inflated sense of winning frequency \citep{dixon2010}. This approach exploits cognitive mechanisms whereby salient feedback enhances the effectiveness of rewards in driving learning and decision-making.

Experimental evidence directly supports this design strategy. Research by \citet{mereu2010} found that visually enhancing win outcomes significantly increased participants' illusion of control—their perception of influence over random outcomes—while similarly enhanced loss feedback had minimal effect. Indeed, research suggests that EGMs with enhanced audiovisual feedback—and especially LDWs—promote persistent gambling even in the face of net losses presumably by increasing play desirability and misperceptions of wins, though effects vary by individual \citep{dixon2014, graydon2018, myles2023}. This suggests that SRCs may not only intensify the experience of genuine rewards but potentially distort the very categorization of outcomes in decision-making processes.

This strategic deployment of salience has intensified with the technological evolution of EGMs, which now incorporate highly sophisticated audiovisual features synchronized with game events \citep{schull2014}. Comparative studies demonstrate that machines with enhanced sensory features, such as richer sounds, are often preferred by players and generate higher revenue, even when the underlying probabilities and payout structures are identical \citep{dixon2014, graydon2018}. In a controlled laboratory experiment, \citet{cherkasova2018} demonstrated that SRCs directly promoted riskier choices in a gambling task, independent of the options' expected values, providing evidence for how such features might influence decision-making. In practical terms, participants essentially chose to play against the mathematical odds simply due to the enhanced sensory experience of winning.

Furthermore, physiological measurements indicate that SRCs elicit stronger autonomic responses, including increased skin conductance and heart rate deceleration—markers of heightened attention and emotional arousal \citep{dixon2014, detez2019}. Understanding the precise cognitive and neural mechanisms driving these effects is therefore crucial.

\subsection{Neurocognitive Mechanisms of Salience and Reward Processing}

Exploring the neurocognitive underpinnings reveals how salience might interact with reward processing. From a cognitive neuroscience perspective, this phenomenon involves multiple brain systems operating in parallel. Research indicates that the brain processes outcome information through distinct yet interacting pathways—an early salience detection system involving regions such as the anterior insula and dorsal anterior cingulate cortex, which are critical for identifying behaviorally relevant stimuli \citep{seeley2007}, followed by value computation in the ventral striatum and ventromedial prefrontal cortex \citep{bartra2013}.

Human fMRI work shows that when monetary rewards are rendered behaviorally salient, activity in the ventral striatum is amplified \citep{zink2004}. At the cellular level, a subset of mid‑brain dopaminergic neurons displays stronger phasic bursts when rewards (or other outcomes) are preceded by novel or attention‑capturing cues, indicating that salience can boost dopaminergic signaling beyond classical value coding, potentially enhancing the effective learning rate beyond what reward magnitude alone would predict \citep{bromberg2010}. Consistent with this, fibre‑photometry recordings reveal larger nucleus‑accumbens dopamine transients as the physical intensity—and hence perceived salience—of appetitive or aversive stimuli increases \citep{kutlu2021}. Collectively, these findings point towards plausible neural mechanisms for the proposed $\alpha$ shift: salient feedback may amplify dopaminergic prediction error signals, leading to a larger update step and thus a higher effective learning rate specifically for saliently signaled outcomes.

Beyond direct modulation of reward signals, SRCs also influence attention and arousal. In the aforementioned study by \citet{cherkasova2018}, eye-tracking revealed that participants spent less time examining probability information when salient win cues were present, suggesting a shift in attentional priorities away from deliberative value assessment. Concurrently, pupillometry indicated greater pupil dilation following wins in the cued condition. This pattern, linked to heightened locus coeruleus-norepinephrine activity, suggests increased physiological arousal elicited by the salient rewards. Together, these findings indicate that sensory salience can shift both what information is prioritized and how strongly rewarding outcomes shape subsequent choices—offering a plausible mechanism through which salience may modulate learning.

\subsection{Individual Differences in Susceptibility to Salient Cues}

An important consideration is that the impact of SRCs likely varies across individuals. Personality traits such as sensation seeking—characterized by the desire for novel, varied, and intense sensory experiences—and impulsivity—the tendency toward rapid, unplanned actions with diminished regard for negative consequences—have been linked to both gambling behavior and heightened neural responses to reward‑predictive cues \citep{hawes2017, verdejo2008, zuckerman1994}. These traits may moderate the effect of salient feedback on learning rates and decision strategies, potentially explaining why some individuals are more vulnerable to environments rich in such cues.

Indeed, higher sensation seeking is associated with stronger psychophysiological and neural responses to reward cues \citep{hawes2017} and potentially biased learning, favoring positive over negative outcomes \citep{xu2019}. Trait impulsivity has also been linked to altered RL profiles, sometimes manifesting as reduced learning from negative feedback compared to rewards \citep{pauli2025, zou2022}. Clinical data also confirm that higher impulsivity co‑occurs with greater gambling severity \citep{blaszczynski1997, slutske2005}. Therefore, these individual differences could interact with environmental factors like feedback salience, potentially amplifying or dampening its effect on learning and contributing to varied susceptibility to maladaptive decision-making in stimulating environments.

\subsection{Research Gaps and Current Study}

While the research reviewed above highlights the potent influence of salient sensory feedback, particularly in contexts like EGM play, several gaps remain in understanding the precise mechanisms through which these cues modulate learning and decision-making. Specifically, further research is needed to quantify these effects within a formal learning framework and to isolate the impact of salience itself.

First, although behavioral consequences like increased gambling persistence are well-documented \citep{dixon2014, graydon2019}, fewer studies have employed computational modeling, such as reinforcement learning, to formally quantify how salient feedback specifically alters underlying learning parameters like the learning rate ($\alpha$). Understanding whether salience directly modifies the weight given to prediction errors is crucial for mechanistic insight. Furthermore, existing research has typically compared machines with different levels of overall sensory stimulation rather than systematically manipulating the presence of salient feedback while holding reward values constant. This methodological limitation has made it difficult to isolate the specific effect of sensory salience on learning processes.

Second, the choice of task is critical. While effects have been observed in various paradigms \citep{cherkasova2018, mereu2010}, tasks specifically designed to probe trial-by-trial adjustments based on feedback, such as multi-armed bandit tasks, are particularly well-suited for quantifying learning rate dynamics using RL models. The use of a restless bandit task, where reward probabilities change over time, further necessitates continuous learning and adaptation, making it sensitive to factors modulating learning speed.

Third, while the general impact of salience on attention and choice is studied, the specific way SRCs—a hallmark of EGM design—influence reinforcement learning parameters and simple behavioral strategies like win-stay behavior requires more targeted investigation in controlled experimental settings.

Finally, although individual differences in traits like sensation seeking and impulsivity are known moderators of reward sensitivity and gambling behavior \citep{hawes2017, pauli2025}, how these traits specifically relate to the magnitude of learning modulation induced by salient feedback (i.e., the $\alpha$ shift) remains an open area for investigation.

The current study is designed to address these gaps using a controlled experimental approach. Specifically, a restless two-armed bandit task is employed where participants make choices under conditions of changing reward probabilities. Crucially, within this task, the presence of salient audiovisual feedback accompanying win outcomes is systematically manipulated using a variable ratio schedule, while the reward outcomes themselves remain independent of the feedback type. This within-subject design allows for direct comparison of behavior and learning parameters between trials with salient versus non-salient feedback for identical win outcomes, thereby isolating the effect of sensory salience.

The primary objective is to investigate how salient, win-paired audiovisual cues modulate learning rates and choice strategies within an RL framework. To quantify these effects, a hierarchical Bayesian reinforcement learning model will be fitted to the trial-by-trial choice data. It is hypothesized that:

\begin{enumerate}
\item The presence of salient feedback following wins will enhance learning from those outcomes, manifesting as a positive $\alpha_{\text{shift}}$ parameter within the RL model. This parameter explicitly captures the difference in effective learning rates between salient and non-salient win trials.

\item Salient feedback following a win will increase the likelihood of repeating the rewarded choice on the subsequent trial, resulting in higher win-stay probability compared to wins followed by non-salient feedback.
\end{enumerate}

Furthermore, the study will explore potential associations between individual differences in sensation seeking and impulsivity and the magnitude of the observed effects of salient feedback on learning ($\alpha_{\text{shift}}$) and choice behavior (win-stay probability).

By integrating computational modeling with a precisely controlled experimental design targeting dynamic learning, this research aims to provide a clearer understanding of how salient environmental cues, inspired by real-world applications like EGMs, can shape adaptive learning processes and potentially contribute to maladaptive choice patterns. 